\subsection{公理}

一个等号理论是一个这样的理论 \(\mathcal{C}\) ：它有一个权重为2的关系符号，写作 =（读作“等于”），并且其中方案S1至S5（§§3和4）连同下面的方案S6和S7一起提供隐式公理。如果 \(T\) 和 \(U\) 是 \(\mathcal{C}\) 中的项，则组装 = \(TU\) 依据 CF4 是 \(\mathcal{C}\) 中的一个关系（称为相等关系）；在实践中它被记为 \(T = U\) 或 \((T) = (U)\) .

S6. 设 x 为一个字母，设 T 和 U 为 C 中的项，且设 \(R\{x\}\) 是 C 中的一个关系；则关系 \((T = U) \Longrightarrow (R\{T\} \Longleftrightarrow R\{U\})\) 是一条公理。

S7. 若R和S是C中的关系，且x是一个字母，则关系 \(((\forall x)(R \leftrightarrow S)) \Longrightarrow (\tau_{x}(R) = \tau_{x}(S))\) 是一条公理。

为了证明规则S6是一个方案，设A是G的一个公理，通过应用S6得到；则G中存在一个关系R，G中的项T和U，以及一个字母x，使得A为 \((T=U)\Longrightarrow((T|x)R\Longleftrightarrow(U|x)R)\) 。我们将证明，如果y是一个字母且V是C中的一个项，则关系 \((V|y)A\) 可以通过应用S6得到。借助CS1（§1，第2节），我们可以假设x与y不同且不出现在V中。设 \(T'\) , \(U'\) , \(R'\) 分别表示组装 \((V|y)T\) , \((V|y)U\) , \((V|y)R\) 。由CS2和CS5（§1，第2节）， \((V|y)A\) 与下式相同：

\[
(T ^ {\prime} = U ^ {\prime}) \Longrightarrow ((T ^ {\prime} | x) R ^ {\prime} \Longleftrightarrow (U ^ {\prime} | x) R ^ {\prime})
\]

证明便完成了。验证S7是一个方案的过程类似。

直观上，方案S6意味着如果两个对象相等，则它们具有相同的性质。方案S7与日常直觉相距较远；它意味着如果对象x的两个性质R和S等价，则被区分的对象 \(\tau_{\mathbf{x}}(\mathbf{R})\) 和 \(\tau_{\mathbf{x}}(\mathbf{S})\) （分别从满足R的对象和满足S的对象中选取，如果这样的对象存在的话）相等。读者会注意到，S7中量词 \(\forall x\) 的存在是本质的（参见练习7）。

关系 = TU 的否定记为 \(T \neq U\) 或 \((T) \neq (U)\) （其中符号 \(\neq\) 读作“不同于”）。

由S6我们推出以下判据：

C43. 设 x 为一个字母，设 T 和 U 为 G 中的项，并设 \(R\{x\}\) 为 G 中的一个关系。则关系

\[
(\boldsymbol {T} = \boldsymbol {U} \text {   且   } \boldsymbol {R} \{\boldsymbol {T} \}), \quad (\boldsymbol {T} = \boldsymbol {U} \text {   且   } \boldsymbol {R} \{\boldsymbol {U} \})
\]

是等价的。

因为如果我们附加假设 \(T = U\) 和 \(R\{T\}\) ，那么 \(R\{U\}\) 由S6为真；因此 \((T = U \text{ 且 } R\{U\})\) 为真。

当一个形如 \(T = U\) 的关系在理论 \(\mathfrak{C}\) 中已被证明时，通常（按语言习惯）说 \(T\) 和 \(U\) “是相同的”或“同一的”。同样地，当 \(T \neq U\) 在 \(\mathfrak{C}\) 中为真时，我们说 \(T\) 和 \(U\) “互不相同”，以代替说 \(T\) 不同于 \(U\) .

